\section*{अध्याय \ref{ch:g2:drawing-to-map} --- चित्र से मानचित्र तक}

\begin{solution}{exo:g2:drawing-to-map:1}
ऊपर से देखे हुए रूप में।
\end{solution}

\begin{solution}{exo:g2:drawing-to-map:2}
पश्चिम की दीवार पर नीली रेखाएँ खिड़कियाँ हैं। उत्तर-पूर्वी कोने में
अलमारी है।
\end{solution}

\begin{solution}{exo:g2:drawing-to-map:3}
एक शीर्षक, एक \omterm{def:g2:drawing-to-map:legend}{संकेत-सूची}, एक \omterm{def:g2:drawing-to-map:scale}{मापनी} और एक \omterm{def:g2:drawing-to-map:directions}{दिशा-चक्र} (या उत्तर की ओर इशारा करता एक तीर)।
\end{solution}

\begin{solution}{exo:g2:drawing-to-map:4}
उत्तर सबसे ऊपर है; पश्चिम बाईं ओर।
\end{solution}

\begin{solution}{exo:g2:drawing-to-map:5}
एक वर्ग 1 मीटर के बराबर है: मेज़ 2 मीटर लंबी और 1 मीटर चौड़ी है।
\end{solution}

\begin{solution}{exo:g2:drawing-to-map:6}
भूरा वर्ग फ़ार्म है; हरी आकृति जंगल है।
\end{solution}

\begin{solution}{exo:g2:drawing-to-map:7}
जंगल \omterm{def:g1:landscapes-around-us:water}{नदी} के दक्षिण में है, और फ़ार्म के पूर्व में।
\end{solution}

\begin{solution}{exo:g2:drawing-to-map:8}
एक वर्ग 1 किलोमीटर के बराबर है: पश्चिम से पूर्व तक 10 किलोमीटर और
दक्षिण से उत्तर तक 8 किलोमीटर।
\end{solution}

\begin{solution}{exo:g2:drawing-to-map:9}
तीन-तीन मेज़ों की तीन पंक्तियाँ: 9 मेज़ें। दरवाज़ा पूर्व की दीवार के दक्षिणी
हिस्से में है, इसलिए दरवाज़े के सबसे पास वाली मेज़ दक्षिण में है (कमरे के
दक्षिण-पूर्वी कोने में)।
\end{solution}

\begin{solution}{exo:g2:drawing-to-map:10}
पलंग 2 वर्ग लंबा है। कमरा 4 वर्ग गुणा 3 वर्ग है: 4-4 वर्गों की 3 पंक्तियाँ,
यानी $4 + 4 + 4 = 12$ वर्ग।
\end{solution}

\begin{solution}{exo:g2:drawing-to-map:11}
``बाईं ओर'' इस पर निर्भर करता है कि लीना का मुँह किधर है: अगर वह घूम जाए,
तो \omterm{def:g1:landscapes-around-us:water}{नदी} उसकी दाईं ओर होगी। सबके लिए सच: ``\omterm{def:g1:landscapes-around-us:water}{नदी} फ़ार्म के उत्तर में है।''
\end{solution}
